\section{Kern und Bild}

In den Übungen zum vorangegangenen Kapitel sind wir auf die folgenden Gleichungssysteme gestossen:

\[
\begin{cases}
-4x + 6y = 5 \\
6x - 9y = -8
\end{cases}
\quad \text{und} \quad
\begin{cases}
-4x + 6y = 5 \\
6x - 9y = -7.5
\end{cases}
\]

Da 
\[
\det \begin{pmatrix} -4 & 6 \\ 6 & -9 \end{pmatrix} = 0,
\]
kann man diese Gleichungssysteme nicht mit der Inversen der Matrix lösen.

Wenn man sie auf konventionelle Art löst, stellt man leicht fest, dass das linke Gleichungssystem keine Lösung hat, das rechte hingegen unendlich viele.

\medskip

In diesem Kapitel werden wir Begriffe und Ideen einführen, mit denen wir einfach feststellen können, ob so ein System lösbar ist.

\medskip

Zunächst einmal wollen wir die Situation geometrisch auffassen:  
Gesucht ist ein Vektor
\[
\vec{v} = \begin{pmatrix} x \\ y \end{pmatrix}
\]
mit der Eigenschaft, dass das Bild \( A\vec{v} \) unter der Matrix \( A \) gerade der Vektor
\[
\begin{pmatrix} 5 \\ -8 \end{pmatrix} \quad \text{bzw.} \quad \begin{pmatrix} 5 \\ -7.5 \end{pmatrix}
\]
ist.

\medskip

Dazu führen wir den Begriff der Bildmenge einer linearen Abbildung ein:
\begin{definition}
    {Bild einer Abbildung}{}
   

Sei \( f \colon \mathbb{R}^2 \to \mathbb{R}^2 \) eine Abbildung.  
Die Menge aller Bilder von \( f \) nennt man das \emph{Bild} von \( f \) und schreibt dafür:
\[
\operatorname{im}(f) = \left\{ f(\vec{v}) \mid \vec{v} \in \mathbb{R}^2 \right\}
\]
Die Bezeichnung \( \operatorname{im} \) steht für \emph{image} (englisch für Bild).

% \medskip

% Die Menge aller Vektoren, die auf den Nullvektor abgebildet werden, nennt man den \emph{Kern} von \( f \) und schreibt dafür:
% \[
% \ker(f) = \left\{ \vec{v} \in \mathbb{R}^2 \mid f(\vec{v}) = \vec{0} \right\}
% \]

\end{definition}
\textbf{Bemerkung}: Manchmal spricht man auch vom Bild einer Matrix und meint dabei das Bild der zugehörigen linearen Abbildung.

Wie bestimmt man konkret das Bild einer linearen Abbildung? Die folgenden Überlegungen geben eine Antwort darauf:


% Die lineare Abbildung \( f \) habe die Matrix
% \[
% A = \begin{pmatrix} -4 & 6 \\ 6 & -9 \end{pmatrix}.
% \]
% Für das Bild eines allgemeinen Vektors
% \[
% \begin{pmatrix} x \\ y \end{pmatrix}
% \]
% gilt:
% \begin{align*}
% A \begin{pmatrix} x \\ y \end{pmatrix}
% &= A \left( x \begin{pmatrix} 1 \\ 0 \end{pmatrix} + y \begin{pmatrix} 0 \\ 1 \end{pmatrix} \right) \\
% &= x \cdot A \begin{pmatrix} 1 \\ 0 \end{pmatrix} + y \cdot A \begin{pmatrix} 0 \\ 1 \end{pmatrix} \\
% &= x \begin{pmatrix} -4 \\ 6 \end{pmatrix} + y \begin{pmatrix} 6 \\ -9 \end{pmatrix}
% \end{align*}

% Anders gesagt:  
% Alle möglichen Bildpunkte von \( f \) können als eine \emph{Linearkombination} der beiden Spaltenvektoren
% \[
% \begin{pmatrix} -4 \\ 6 \end{pmatrix} \quad \text{und} \quad \begin{pmatrix} 6 \\ -9 \end{pmatrix}
% \]
% dargestellt werden.

% \medskip

% In diesem speziellen Beispiel kann das Resultat noch vereinfacht werden, wenn man berücksichtigt, dass die Spalten \emph{linear abhängig} sind:
% \[
% \begin{pmatrix} 6 \\ -9 \end{pmatrix} = -\frac{3}{2} \cdot \begin{pmatrix} -4 \\ 6 \end{pmatrix}.
% \]

% Damit folgt:



% \begin{align*}
% A \begin{pmatrix} x \\ y \end{pmatrix}
% &= x \begin{pmatrix} -4 \\ 6 \end{pmatrix} + y \begin{pmatrix} 6 \\ -9 \end{pmatrix} \\
% &= x \begin{pmatrix} -4 \\ 6 \end{pmatrix} - \frac{3}{2} y \begin{pmatrix} -4 \\ 6 \end{pmatrix} \\
% &= \underbrace{\left(x - \frac{3}{2} y\right)}_{=\,t} \begin{pmatrix} -4 \\ 6 \end{pmatrix} \\
% &= t \begin{pmatrix} -4 \\ 6 \end{pmatrix}
% \end{align*}

% Die Vektoren in \( \operatorname{im}(f) \) sind demnach alle Vielfache von \( \begin{pmatrix} -4 \\ 6 \end{pmatrix} \).  
% Somit gilt:
% \[
% \operatorname{im}(f) = \left\{ t \begin{pmatrix} -4 \\ 6 \end{pmatrix} \;\middle|\; t \in \mathbb{R} \right\}.
% \]

% Anders gesagt:  
% Unabhängig davon, welchen Vektor \( \vec{v} \in \mathbb{R}^2 \) man wählt, liegt \( A \cdot \vec{v} \) sicher auf der Geraden mit der Parametergleichung
% \[
% \vec{r}(t) = t \begin{pmatrix} -4 \\ 6 \end{pmatrix}.
% \]

% Das erklärt auch, warum das Gleichungssystem
% \[
% \begin{cases}
% -4x + 6y = 5 \\
% 6x - 9y = -8
% \end{cases}
% \]
% \emph{keine Lösung} hat: Der rechte Vektor \( \begin{pmatrix} 5 \\ -8 \end{pmatrix} \) liegt nicht auf der Geraden \( \vec{r}(t) \)!

% \medskip

% \textbf{Kern:} Der Kern von \( f \) ist die Menge aller Vektoren, die auf den Nullvektor abgebildet werden.  
% Mit anderen Worten: Der Kern von \( f \) ist die Lösungsmenge der Gleichung
% \[
% A \begin{pmatrix} x \\ y \end{pmatrix} = \vec{0}.
% \]

% Analog zur obigen Rechnung ergibt sich:
% \begin{align*}
% A \begin{pmatrix} x \\ y \end{pmatrix}
% &= x \begin{pmatrix} -4 \\ 6 \end{pmatrix} + y \begin{pmatrix} 6 \\ -9 \end{pmatrix} \\
% &= \left(x - \frac{3}{2} y\right) \begin{pmatrix} -4 \\ 6 \end{pmatrix} = \vec{0}
% \end{align*}

% Daraus folgt:
% \[
% x - \frac{3}{2} y = 0 \quad \Rightarrow \quad x = \frac{3}{2} y
% \quad \Rightarrow \quad
% \begin{pmatrix} x \\ y \end{pmatrix} = t \begin{pmatrix} 3 \\ 2 \end{pmatrix}, \quad t \in \mathbb{R}.
% \]

% Also:
% \[
% \ker(f) = \left\{ t \begin{pmatrix} 3 \\ 2 \end{pmatrix} \;\middle|\; t \in \mathbb{R} \right\}.
% \]

% Anders gesagt: Für alle Vektoren \( \vec{w} \), welche auf der Geraden
% \[
% \vec{r}(t) = t \begin{pmatrix} 3 \\ 2 \end{pmatrix}
% \]
% liegen, gilt:
% \[
% A \cdot \vec{w} = \vec{0}.
% \]

\begin{bsp}
    { Bild von  \( A = \begin{pmatrix} -4 & 6 \\ 6 & -9 \end{pmatrix} \) und Lösbarkeit der Gleichungssysteme}{}

    \vspace*{19 cm}
\end{bsp}
\clearpage
\begin{satz}
    {Spalten und Bild}{}
    
Seien \( \vec{v}_1 \) und \( \vec{v}_2 \) die Spalten einer \( 2 \times 2 \)-Matrix \( A \).

\medskip

Das \textbf{Bild} der Matrix ist die Menge aller Linearkombinationen der Spaltenvektoren:
\[
\operatorname{im}(A) = \left\{ t \cdot \vec{v}_1 + s \cdot \vec{v}_2 \;\middle|\; t, s \in \mathbb{R} \right\}
\]

\medskip
Ein Gleichungssystem der Form 
$$ A\cdot \vec{v} = \vec{b}, \qquad \vec{v} =? $$
ist genau dann lösbar, wenn der Vektor \( \vec{b} \in \operatorname{im}(A)  \).

% \medskip

% Der \textbf{Kern} einer linearen Abbildung besteht aus den Vektoren \( \begin{pmatrix} x \\ y \end{pmatrix} \), für welche gilt:
% \[
% x \cdot \vec{v}_1 + y \cdot \vec{v}_2 = \vec{0}
% \]

\end{satz}


\begin{bsp}
{}{}
Das Gleichungssystem $\begin{cases}
-4x + 6y = 5 \\
6x - 9y = -8
\end{cases}$ ist nicht lösbar, weil der Vektor \( \begin{pmatrix} 5 \\ -8 \end{pmatrix} \) nicht  im Bild $\operatorname{im}(A) = \left\{ t\cdot \begin{pmatrix}
    2\\-3 
\end{pmatrix} \mid t \in \mathbb{R} \right\}$ der Matrix $A=\begin{pmatrix}
-4 & 6 \\ 6 & -9 \end{pmatrix}$ liegt.

\end{bsp}



Das zweite Gleichungssystem
\[\begin{cases}
-4x + 6y = 5 \\
6x - 9y = -7.5
\end{cases}\]
hat hingegen unendlich viele Lösungen. Diese haben die Form
\[ \begin{pmatrix}
x \\ y \end{pmatrix} = \begin{pmatrix}
    5 \\ -\frac{5}{2}
\end{pmatrix} + t \begin{pmatrix} 3 \\ 2 \end{pmatrix}, \quad t \in \mathbb{R}.\]








































\begin{bsp}
    {Kern und Bild von $A = \begin{pmatrix} 1 & 3 \\ 3 & 9 \end{pmatrix}$ }{}

% Gegeben sei die Matrix
% \[
% A = \begin{pmatrix} 1 & 3 \\ 3 & 9 \end{pmatrix}.
% \]

Die beiden Spalten sind linear abhängig.  
Also gilt:
\[
\operatorname{im}(A) = \left\{ t \begin{pmatrix} 1 \\ 3 \end{pmatrix} \;\middle|\; t \in \mathbb{R} \right\}.
\]

Dreimal die erste Spalte minus die zweite ergibt den Nullvektor, also ist:
\[
\ker(A) = \left\{ t \begin{pmatrix} 3 \\ -1 \end{pmatrix} \;\middle|\; t \in \mathbb{R} \right\}.
\]

\end{bsp}

\begin{bsp}
{Kern und Bild von $B = \begin{pmatrix} 1 & 3 \\ 4 & -1 \end{pmatrix}.$}{}

% Gegeben sei die Matrix
% \[
% B = \begin{pmatrix} 1 & 3 \\ 4 & -1 \end{pmatrix}.
% \]

Die beiden Spalten sind linear unabhängig.  
Also gilt:
\[
\operatorname{im}(B) = \left\{ t \begin{pmatrix} 1 \\ 4 \end{pmatrix} + s \begin{pmatrix} 3 \\ -1 \end{pmatrix} \;\middle|\; t, s \in \mathbb{R} \right\} = \mathbb{R}^2.
\]

Da die Spalten linear unabhängig sind, hat die Gleichung
\[
x \cdot \vec{v}_1 + y \cdot \vec{v}_2 = \vec{0}
\]
nur die triviale Lösung \( x = 0 \), \( y = 0 \).  
Also gilt:
\[
\ker(B) = \left\{ \vec{0} \right\}.
\]

\end{bsp}



Die beiden Beispiele machen deutlich, dass Bild und der Kern einer linearen Abbildung ganz spezielle Teilmengen von $\R^2$ bzw. $\R^3$ sind. Es sind sogenannte \emph{Unterräume}:













\pagebreak
\begin{definition}
    {Unterraum}{}

Ein \textbf{Unterraum} von \( \mathbb{R}^2 \)  ist eine Teilmenge \( U \subseteq \mathbb{R}^2 \) mit folgenden Eigenschaften:
\begin{itemize}
  \item Für alle \( \vec{v}, \vec{w} \in U \) gilt: \( \vec{v} + \vec{w} \in U \)
  \item Für alle \( \vec{v} \in U \) und alle \( \lambda \in \mathbb{R} \) gilt: \( \lambda \vec{v} \in U \)
\end{itemize}

    \end{definition}

\textbf{Bemerkungen:} 

\begin{itemize}
  \item Ein Unterraum enthält immer den Nullvektor: \( \vec{0} \in U \).
  
  \item Es gibt drei Arten von Unterräumen in \( \mathbb{R}^2 \):
  \begin{itemize}
    \item Der \textbf{nulldimensionale} Unterraum:  
    \[
    U = \{ \vec{0} \}
    \]
    
    \item \textbf{Eindimensionale} Unterräume der Form:
    \[
    U = \{ t \cdot \vec{v} \mid t \in \mathbb{R} \}
    \]
    Dies sind genau die Geraden durch den Ursprung.

    \item Der \textbf{zwei-dimensionale} Unterraum von \( \mathbb{R}^2 \):
    \[
    U = \mathbb{R}^2
    \]
  \end{itemize}
\end{itemize}
\begin{satz}
    {}{}
    Kern und Bild einer linearen Abbildung sind Unterräume von \( \mathbb{R}^2 \).
\end{satz}
\begin{proof}
    
Wir beweisen die Aussage für den Kern einer Matrix bzw. linearen Abbildung.

\begin{itemize}
    \item Seien \( \vec{v} \) und \( \vec{w} \) zwei Vektoren in \( \ker(f) \).  
Weil \( f \) linear ist, gilt:
\[
f(\vec{v} + \vec{w}) = f(\vec{v}) + f(\vec{w}) = \vec{0} + \vec{0} = \vec{0}
\]
Somit ist \( \vec{v} + \vec{w} \in \ker(f) \).

\item Sei nun \( \vec{v} \in \ker(f) \) und \( \lambda \in \mathbb{R} \).  
Wegen der Linearität von \( f \) gilt:
\[
f(\lambda \vec{v}) = \lambda f(\vec{v}) = \lambda \cdot \vec{0} = \vec{0}
\]
Also ist auch \( \lambda \vec{v} \in \ker(f) \).
\end{itemize}


\end{proof}


\begin{satz}
    {Äquivalente Aussagen für lineare Abbildungen}{}

Sei \( f \colon \mathbb{R}^2 \to \mathbb{R}^2 \) eine lineare Abbildung mit der zugehörigen Matrix \( A \).  
Dann sind die folgenden Aussagen äquivalent:

\begin{enumerate}
  \item \( \det(A) \neq 0 \)
  \item Die Spalten von \( A \) sind linear unabhängig.
  \item \( \operatorname{im}(f) = \mathbb{R}^2 \)
  \item \( \ker(f) = \{ \vec{0} \} \)
\end{enumerate}

\end{satz}
\begin{proof}
    

Wir beweisen die folgenden Implikationen:  
\[
1 \Rightarrow 2 \Rightarrow 3 \Rightarrow 4 \Rightarrow 1
\]

\begin{description}
  \item[\(1 \Rightarrow 2\): (Beweis durch Widerspruch)]
  Angenommen, es gäbe eine Matrix \( A \) mit \( \det(A) \neq 0 \), deren Spalten linear abhängig sind.  
  Dann hat \( A \) die Form:
  \[
  A = \begin{pmatrix} a & k \cdot a \\ b & k \cdot b \end{pmatrix}
  \]
  Damit gilt:
  \[
  \det(A) = a \cdot k b - k a \cdot b = 0
  \]
  Widerspruch zu \( \det(A) \neq 0 \). Also müssen die Spalten linear unabhängig sein.

  \item[\(2 \Rightarrow 3\):]  
  Das folgt direkt daraus, dass das Bild einer linearen Abbildung aus den Linearkombinationen der Spalten der Matrix besteht.  
  Sind die Spalten linear unabhängig, so spannen sie den gesamten Raum \( \mathbb{R}^2 \) auf. Also:
  \[
  \operatorname{im}(f) = \mathbb{R}^2.
  \]

  \item[\(3 \Rightarrow 4\):]  
  Seien \( \vec{v}_1 \) und \( \vec{v}_2 \) die beiden Spaltenvektoren der Matrix.  
  Da das Bild von \( f \) aus den Linearkombinationen der Spalten besteht und \( \operatorname{im}(f) = \mathbb{R}^2 \) ist, müssen die Spalten linear unabhängig sein.  
  Damit hat die Gleichung
  \[
  x \cdot \vec{v}_1 + y \cdot \vec{v}_2 = \vec{0}
  \]
  nur die triviale Lösung \( x = y = 0 \), und es folgt:
  \[
  \ker(f) = \{ \vec{0} \}
  \]

  \item[\(4 \Rightarrow 1\): (Beweis durch Widerspruch)]
  Sei:
  \[
  A = \begin{pmatrix} a & b \\ c & d \end{pmatrix}, \quad \text{und angenommen: } \det(A) = 0.
  \]
  Betrachte den Vektor:
  \[
  \vec{v} = \begin{pmatrix} d \\ -c \end{pmatrix}
  \]
  Dann gilt:
  \[
  f(\vec{v}) = A \cdot \vec{v} = 
  \begin{pmatrix} a & b \\ c & d \end{pmatrix}
  \begin{pmatrix} d \\ -c \end{pmatrix}
  = \begin{pmatrix} ad - bc \\ cd - dc \end{pmatrix}
  = \begin{pmatrix} \det(A) \\ 0 \end{pmatrix}
  = \vec{0}
  \]
  Also \( \vec{v} \in \ker(f) \Rightarrow \vec{v} = \vec{0} \Rightarrow \begin{pmatrix} d \\ -c \end{pmatrix} = \begin{pmatrix} 0 \\ 0 \end{pmatrix} \)

  Analog für:
  \[
  \vec{w} = \begin{pmatrix} b \\ -a \end{pmatrix}, \quad
  f(\vec{w}) = A \cdot \vec{w} = \begin{pmatrix} ab - ba \\ cb - da \end{pmatrix}
  = \begin{pmatrix} 0 \\ -\det(A) \end{pmatrix}
  = \vec{0}
  \Rightarrow \vec{w} = \vec{0}
  \]

  Damit hat \( A \) die Form:
  \[
  A = \begin{pmatrix} 0 & 0 \\ 0 & 0 \end{pmatrix}
  \]
  Widerspruch zu \( \ker(f) = \{ \vec{0} \} \).  
  Also muss gelten: \( \det(A) \neq 0 \).
\end{description}

\end{proof}

\begin{ueb}
    {}{}

\begin{enumerate}
  \item Bestimmen Sie den Kern und das Bild der folgenden Matrizen.  
  \begin{enumerate}
    \item \( A = \begin{pmatrix} 1 & 2 \\ 3 & 4 \end{pmatrix} \)
    \item \( B = \begin{pmatrix} 0 & 0 \\ 0 & 0 \end{pmatrix} \)
    \item \( C = \begin{pmatrix} 2 & 12 \\ -1 & -6 \end{pmatrix} \)
    \item \( D = \begin{pmatrix} -6 & 7 \\ 12 & -14 \end{pmatrix} \)
    \item \( E = \begin{pmatrix} 1 & 3 \\ 0 & 0 \end{pmatrix} \)
  \end{enumerate}

  \item Sei \( f \) die Projektion auf die Gerade \( y = x \) in Richtung  
  \( \vec{v} = \begin{pmatrix} 1 \\ -1 \end{pmatrix} \).  
  Bestimmen Sie den Kern und das Bild von \( f \).

  \item Geben Sie eine Matrix an, deren Kern die Gerade  
  \( g \colon \vec{r}(t) = t \begin{pmatrix} 1 \\ 2 \end{pmatrix} \)  
  und deren Bild die \( y \)-Achse ist.

  \item Geben Sie eine Matrix \( A \) an, deren Kern und Bild übereinstimmen.  
  Was ergibt \( A^2 \)?

  \item Sei \( f \) die Spiegelung an der \( xz \)-Ebene.  
  Bestimmen Sie den Kern und das Bild.

  \item Sei \( f \) die orthogonale Projektion auf die \( xz \)-Ebene.  
  Bestimmen Sie den Kern und das Bild.

  \item \textbf{Challenges:}
  \begin{enumerate}
    \item Beweisen Sie, dass das Bild einer linearen Abbildung ein Unterraum ist.\\
    \emph{Tipp:} \( \vec{w} \) ist genau dann ein Element von \( \operatorname{im}(f) \), wenn es ein  
    \( \vec{v} \in \mathbb{R}^2 \) mit \( f(\vec{v}) = \vec{w} \) gibt. Gibt es dann einen Vektor  
    \( \vec{x} \) mit \( f(\vec{x}) = \lambda \vec{w} \)? Welchen?

    \item Beweisen Sie die folgende Aussage:  
    Sei \( f \) eine lineare Abbildung mit \( \ker(f) = \{ \vec{0} \} \).  
    Dann ist \( f \) injektiv.\\
    \emph{Tipp:} Injektiv bedeutet: Aus \( f(\vec{v}) = f(\vec{w}) \) folgt \( \vec{v} = \vec{w} \).  
    Betrachten Sie \( f(\vec{v} - \vec{w}) \).
  \end{enumerate}
\end{enumerate}

\end{ueb}
\clearpage
\emph{Lösungen:}
